%%%%%%%%%%%%%%%%%%%%%%%%%%%%%%%%%%%%%%%%%%%%%%%%%%%%%%%%%%%%%%%%%%%%%%%%%%%%%%%%
%                                                                              %
% BlackArch Linux Guide                                                        %
%                                                                              %
%%%%%%%%%%%%%%%%%%%%%%%%%%%%%%%%%%%%%%%%%%%%%%%%%%%%%%%%%%%%%%%%%%%%%%%%%%%%%%%%

\documentclass[a4paper, oneside, 11pt]{book}

%%% INCLUDES %%%
\renewcommand{\familydefault}{\sfdefault}

\usepackage{array}
\usepackage{color}
\usepackage{enumerate}
\usepackage{fancyhdr}
\usepackage{fancyvrb}
\usepackage{geometry}
\usepackage{graphicx}
\usepackage{html}
\usepackage{hyperref}
\usepackage{ifpdf}
\usepackage{listings}
\usepackage{pstricks}
\usepackage{supertabular}
\usepackage{tocloft}
\usepackage[utf8]{inputenc}

%%% LAYOUT %%%
\setlength{\parindent}{0em}
\setlength{\parskip}{1.5ex plus0.5ex minus0.5ex}
\geometry{left=2.5cm, textwidth=16cm, top=3cm, textheight=25cm, bottom=3cm}
\widowpenalty=2000
\clubpenalty=1000
\frenchspacing
\sloppy
\pagecolor[HTML]{FFFFFF}
\color[HTML]{333333}
\setcounter{tocdepth}{10}
\setcounter{secnumdepth}{10}

\hypersetup{
  pdftitle={BlackArch Linux, The BlackArch Linux Guide},
  pdfsubject={BlackArch Linux, The BlackArch Linux Guide},
  pdfauthor={BlackArch Linux, BlackArch Linux},
  pdfkeywords={BlackArch Linux, Penetration Testing, Security, Hacking, Linux},
  pdfcenterwindow=true,
  colorlinks=true,
  breaklinks=true,
  linkcolor=red,
  menucolor=red,
  urlcolor=red
}

% syntax highlighting
% all options should be set here document wide
\lstset{
backgroundcolor=\color[HTML]{EEEEEE},
frame=single,
basicstyle=\footnotesize\ttfamily,
float,
deletekeywords={return},
otherkeywords={mkdir, curl, sudo, sha1sum, grep, cut, sort, wget, makepkg,
pacman, blackman},
keywordstyle=\color{orange},
commentstyle=\color{blue},
stringstyle=\color{red},
language=bash,
showspaces=false,
showtabs=false,
tabsize=2
}

%%% HEADER / FOOTER %%%
\setlength{\headheight}{33pt}
\setlength{\headsep}{33pt}
\lhead{{\includegraphics[width=1cm,height=1cm]{images/logo.png}}}
\rhead{The BlackArch Linux Guide}

%%% CUSTOM MACROS %%%
% for html links
\ifpdf\else
\def\href#1#2{\htmladdnormallink{#2}{#1}}
\fi

%------------------%
%  TITLE PAGE      %
%------------------%
\begin{document}
\pagestyle{empty}
\begin{center}
\begin{figure}[htbp]
\centering
\vspace{0.5cm}
\includegraphics[width=8cm]{images/logo.png}
\label{fig:logo}
\end{figure}
\vspace{0.5cm}
\Huge{\textbf{BlackArch指南}}\\
\vspace{1cm}
\Large{\color{red}https://www.blackarch.org/}\\
\vspace{0.5cm}
\end{center}
\newpage
\tableofcontents
\newpage
\pagestyle{fancy}

%------------------%
%  Chapter 1       %
%------------------%

\chapter{引言}

\section{概览}
BlackArch Linux指南分为以下几部分
\begin{itemize}
\item 引言 - 提供概述，简介和其他有帮助项目的信息
\item 用户指南 - 一个典型用户高效的使用BlackArch所需要知道的一切
\item 开发者指南 - 怎样加入我们，为BlackArch开发和贡献
\item 工具指南 - 工具的详细信息和使用示例
\end{itemize}

\section{什么是BlackArch Linux}
BlackArch是一个为渗透测试工作者和安全研究人员设计的Linux发行版
BlackArch是 \href{https://www.archlinux.org/}{ArchLinux}的一个衍生版，用户可以单独或者成组的安装BlackArch的组件到 \href{https://www.archlinux.org/}{ArchLinux}中
工具集也被发布为一个Arch Linux \href{https://wiki.archlinux.org/index.php/Unofficial\_User\_Repositories}
{unofficial user repository} 这样你可以在一个已经安装好的Arch Linux上面安装BlackArch。 软件包可以独立安装也可以按类别安装。

% should quickly describe the testing methods/code review procedures etc
仓库目前包括超过 \href{https://www.blackarch.org/tools.html}{1300} 种工具，并且工具数目在不断增长。
所以工具在添加到代码库之前都经过仔细的检查以保证仓库的质量。
\section{BlackArch Linux的历史}
敬请期待

\section{支持的平台}
敬请期待

\section{参与贡献}
You can get in touch with the BlackArch team using the following avenues:
你可以通过以下几种途径与BlackArch team取得联系
Website: \url{https://www.blackarch.org/}

Mail: \href{mailto:team@blackarch.org}{team@blackarch.org}

IRC: \url{irc://irc.freenode.net/blackarch}

Twitter: \url{https://twitter.com/blackarchlinux}

Github: \url{https://github.com/Blackarch/}

%------------------%
%  Chapter 2       %
%------------------%


\chapter{用户指南}

\section{安装}

以下章节将会告诉你如何配置BlackArch仓库和安装软件。
BlackArch支持通过二进制包安装也支持从源码编译安装。
BlackArch与普通Arch兼容，其作为一个非官方用户仓库，如果你想要ISO文件，请前往\href{https://www.blackarch.org/downloads.html#iso}{Live ISO}

\subsection{在Arch Linux基础上安装}

用root权限运行 \href{https://blackarch.org/strap.sh}{strap.sh}，跟从指示。示例：
\begin{lstlisting}
   curl -O https://blackarch.org/strap.sh
   sha1sum strap.sh # should match: 86eb4efb68918dbfdd1e22862a48fda20a8145ff
   sudo ./strap.sh
\end{lstlisting}

现在下载最新的软件包列表并且同步软件包:
\begin{lstlisting}
  sudo pacman -Syyu
\end{lstlisting}


\subsection{安装软件包}
你现在可以从blackarch仓库安装工具了
\begin{enumerate}
\item 列出所有可用工具，运行
\begin{lstlisting}
  pacman -Sgg | grep blackarch | cut -d' ' -f2 | sort -u
\end{lstlisting}

\item 安装全部工具，运行:
\begin{lstlisting}
  pacman -S blackarch
\end{lstlisting}

\item 安装一类工具，运行
\begin{lstlisting}
  pacman -S blackarch-<category>
\end{lstlisting}

\item 查看blackarch的目录，运行
\begin{lstlisting}
  pacman -Sg | grep blackarch
\end{lstlisting}

\end{enumerate}

\subsection{从源码安装软件包}
作为一种可选的安装方式，你可以从源码编译安装BlackArch软件包，你可以在\href{https://github.com/BlackArch/blackarch/tree/master/packages}{github}上找到需要的PKGBUILD。
编译整个仓库，可以使用 \href{https://github.com/BlackArch/blackman}{Blackman} 工具。
\begin{itemize}
\item 首先，你需要安装Blackman。如果你的机器上已经配置好了BlackArch包仓库，你可以通过下面的命令安装Blackman:
\begin{lstlisting}
  pacman -S blackman
\end{lstlisting}

\item 你也可以从源码编译安装Blackman:
\begin{lstlisting}
  mkdir blackman
  cd blackman
  wget https://raw2.github.com/BlackArch/blackarch/master/packages/blackman/PKGBUILD
  # Make sure the PKGBUILD has not been maliciously tampered with.
  makepkg -s
\end{lstlisting}

\item 或者从Aur安装blackman:
\begin{lstlisting}
  <whatever AUR helper you use> -S blackman
\end{lstlisting}

\end{itemize}

\subsection{Blackman基本用法} 
尽管用法与pacman不大一样，但是Blackman仍然非常易于使用。下面列出了一些基本用法。
\begin{itemize}
\item 下载，编译安装软件包:
\begin{lstlisting}
  sudo blackman -i package
\end{lstlisting}

\item 下载，编译安装整个类别:
\begin{lstlisting}
  sudo blackman -g group
\end{lstlisting}

\item 下载，编译安装所有BlackArch工具:
\begin{lstlisting}
  sudo blackman -a
\end{lstlisting}

\item 列出BlackArch的类别:
\begin{lstlisting}
  blackman -l
\end{lstlisting}

\item 列出一个类别中的工具:
\begin{lstlisting}
  blackman -p category
\end{lstlisting}

\end{itemize}

\subsection{从live-ISO，netinstall-ISO或者ArchLinux安装}

你可以从live-ISO或者netinstall-ISO中安装BlackArch 查看 \url{https://www.blackarch.org/download.html#iso}。从ISO启动后执行下面的步骤。
\begin{itemize}
\item 安装blackarch-installer:
\begin{lstlisting}
  sudo pacman -S blackarch-installer
\end{lstlisting}

\item 执行
\begin{lstlisting}
  sudo blackarch-install
\end{lstlisting}

\end{itemize}

%------------------%
%  Chapter 3       %
%------------------%

\chapter{开发者指南}

\section{Arch的构建系统和仓库}

PKGBUILD files are build scripts. Each one tells makepkg(1) how to create a
package. PKGBUILD files are written in Bash.
PKGBUILD文件是一种构建脚本。PKGBUILD文件告诉makepkg(1)怎么样去构建一个软件包。PKGBUILD文件是用Bash写的。
For more information, read (or skim through) the following:
想获取更多信息，阅读以下内容：
\begin{itemize}
\item \href{https://wiki.archlinux.org/index.php/Creating_Packages}{Arch Wiki: Creating Packages}
\item \href{https://wiki.archlinux.org/index.php/Makepkg}{Arch Wiki: makepkg}
\item \href{https://wiki.archlinux.org/index.php/PKGBUILD}{Arch Wiki: PKGBUILD}
\item \href{https://wiki.archlinux.org/index.php/Arch_Packaging_Standards}{Arch Wiki: Arch Packaging Standards}
\end{itemize}

\section{Blackarch PKGBUILD 标准}
出于简单，我们的PKGBUILD文件跟AUR的非常相似，下面列出了不同之处。每个软件包至少要属于blackarch，也会有很多软件包属于多个group。
\subsection{Groups}
为了让用户能快速便捷的安装特定需求的一系列软件包，软件包被分为不同的群。群让用户能够简单的通过 "pacman -S <group name>" 来获取一系列软件。
\subsubsection{blackarch}
blackarch群是所有软件包必须属于的基本群。这让用户可以轻而易举的安装所以软件包。
所有软件包都应该属于blackarch。

\subsubsection{blackarch-anti-forensic}
此群中包含用于反取证的软件包，包括加密、隐写，和修改文件属性的软件。包含基本所有以隐藏信息为目的修改系统的软件。

例如: luks, TrueCrypt, Timestomp, dd, ropeadope, secure-delete
\subsubsection{blackarch-automation}
用于工具或工作流自动化的软件包。

例如: blueranger, tiger, wiffy

\subsubsection{blackarch-backdoor}
用于利用或打开漏洞系统上后门的软件包。

例如: backdoor-factory, rrs, weevely

\subsubsection{blackarch-binary}
用于操作二进制文件的软件包。

例如: binwally, packerid

\subsubsection{blackarch-bluetooth}
用于攻击蓝牙标准(802.15.1)的软件包。

例如: ubertooth, tbear, redfang

\subsubsection{blackarch-code-audit}
Packages that audit existing source code for vulnerability analysis.
用于代码审计的软件包。

例如: flawfinder, pscan

\subsubsection{blackarch-cracker}
用于破解密码学函数(比如哈希函数)的软件包。

例如: hashcat, john, crunch

\subsubsection{blackarch-crypto}
除了破解之外的密码学软件包。

例如: ciphertest, xortool, sbd

\subsubsection{blackarch-database}
用于数据库攻击的软件包。

例如: metacoretex, blindsql

\subsubsection{blackarch-debugger}
用于调试程序的软件包。

例如: radare2, shellnoob

\subsubsection{blackarch-decompiler}
用于将编译好的程序逆向转换为源代码的软件包。

例如: flasm, jd-gui

\subsubsection{blackarch-defensive}
用于阻止恶意软件或其他用户攻击的软件包。

例如: arpon, chkrootkit, sniffjoke

\subsubsection{blackarch-disassembler}
This is similar to blackarch-decompiler, and there will probably be a lot
of programs that fall into both, however these packages produce assembly output
rather than the raw source code.
与blackarch-decompiler相似，有许多程序会在两组中同时出现。blackarch-disassembler中的软件会生成汇编代码而不是源代码。

例如: inguma, radare2

\subsubsection{blackarch-dos}
用于Dos(拒绝服务)攻击的软件包。

例如: 42zip, nkiller2

\subsubsection{blackarch-drone}
用于管理无人机的软件包。

例如: meshdeck, skyjack

\subsubsection{blackarch-exploitation}
Packages that takes advantages of exploits in other programs or services.
用于利用其他程序或服务漏洞的软件包。

例如: armitage, metasploit, zarp

\subsubsection{blackarch-fingerprint}
用于指纹识别的软件包。

例如: dns-map, p0f, httprint

\subsubsection{blackarch-firmware}
用于利用固件漏洞的软件包。

例如: 暂无.

\subsubsection{blackarch-forensic}
用于从物理磁盘或者内存中寻找数据的软件包。

例如: aesfix, nfex, wyd

\subsubsection{blackarch-fuzzer}
用于进行fuzz检测(例如对目标进行随机输入,观察结果)的软件包。

例如: msf, mdk3, wfuzz

\subsubsection{blackarch-hardware}
用于攻击或管理物理硬件的软件包。

例如: arduino, smali

\subsubsection{blackarch-honeypot}
用于搭建蜜罐的软件包。

例如: artillery, bluepot, wifi-honey

\subsubsection{blackarch-keylogger}
用于记录键盘的软件包。

例如: 暂无.

\subsubsection{blackarch-malware}
恶意软件以及检测恶意软件的软件包。

例如: malwaredetect, peepdf, yara

\subsubsection{blackarch-misc}
不属于其他分类的软件包。

例如: oh-my-zsh-git, winexe, stompy

\subsubsection{blackarch-mobile}
用于操纵移动平台的软件包。

例如: android-sdk-platform-tools, android-udev-rules

\subsubsection{blackarch-networking}
关于IP网络的软件包。

例如: 基本所有

\subsubsection{blackarch-nfc}
用于nfc(近场通讯)的软件包。

例如: nfcutils

\subsubsection{blackarch-packer}
有关捆绑者的软件包。
\textit{捆绑者是将恶意软件嵌入其他可执行程序中的程序。}

例如: packerid

\subsubsection{blackarch-proxy}
有代理行为的软件包。

例如: burpsuite, ratproxy, sslnuke

\subsubsection{blackarch-recon}
用于主动探测的工具集合。

例如: canri, dnsrecon, netmask

\subsubsection{blackarch-reversing}
反编译器，反汇编器和类似程序的集合。

例如: capstone, radare2, zerowine

\subsubsection{blackarch-scanner}
用于扫描选定系统漏洞的软件包。

Examples: scanssh, tiger, zmap

\subsubsection{blackarch-sniffer}
有关网络流量分析的软件包。

例如: hexinject, pytactle, xspy

\subsubsection{blackarch-social}
主要攻击社交网站的软件包。

例如: jigsaw, websploit

\subsubsection{blackarch-spoof}
用于进行欺骗攻击的软件包。

例如: arpoison, lans, netcommander

\subsubsection{blackarch-threat-model}
用于在特定场景报告或记录威胁模型的软件包。

例如: magictree

\subsubsection{blackarch-tunnel}
用于在给定网络建立隧道的软件包。

例如: ctunnel, iodine, ptunnel

\subsubsection{blackarch-unpacker}
用于从可执行程序中提取恶意软件的软件包。

例如: js-beautify

\subsubsection{blackarch-voip}
用于操作voip程序和协议的软件包。

例如: iaxflood, rtp-flood, teardown

\subsubsection{blackarch-webapp}
用于操作面向互联网的程序的软件包。

例如: metoscan, whatweb, zaproxy

\subsubsection{blackarch-windows}
This group is for any native Windows package that runs via wine.
此群包含所有通过wine运行的Windows程序。

例如: 3proxy-win32, pwdump, winexe

\subsubsection{blackarch-wireless}
用于操作无线网络的软件包。

例如: airpwn, mdk3, wiffy

\section{Repository structure}
你可以在这里找到Blackarch的主要仓库:
\href{https://github.com/BlackArch/blackarch}{https://github.com/BlackArch/blackarch}.
其余次要仓库:
\href{https://github.com/BlackArch}{https://github.com/BlackArch}.

在主git repo中,有三个重要目录:

\begin{itemize}
\item docs - 文档.
\item packages - PKGBUILD 文件.
\item scripts - 有用的小脚本.
\end{itemize}

\subsection{Scripts}
下面是 \verb|scripts/| 目录中的脚本:

\begin{itemize}
\item baaur - Soon, this will upload packages to the AUR.
\item babuild - Build a package.
\item bachroot - Manage a chroot for testing.
\item baclean - Clean old .pkg.tar.xz files from the package repo.
\item baconflict - Soon this will replace \verb|scripts/conflicts|.
\item bad-files - Find bad files in built packages.
\item balock - Obtain or release the package repo lock.
\item banotify - Notify IRC about package pushes.
\item barelease - Release packages to the package repo.
\item baright - Print the BlackArch copyright info.
\item basign - Sign packages.
\item basign-key - Sign a key.
\item blackman - This behaves sort of like pacman but builds from git (not to be
    confused with nrz's Blackman).
\item check-groups - Check groups.
\item checkpkgs - Check packages for errors.
\item conflicts - Check for file conflicts.
\item dbmod - Modify a package database.
\item depth-list - Create a list sorted by dependency depth.
\item deptree - Create a dependency tree, listing only blackarch-provided packages.
\item get-blackarch-deps - Get a list of blackarch dependencies for a package.
\item get-official - Get official packages for release.
\item list-loose-packages - List packages that are not in groups and are not
    dependencies for other packages.
\item list-needed - List missing dependencies.
\item list-removed - List packages that are in the package repo but not in git.
\item list-tools - List tools.
\item outdated - Look for packages that are out-dated in the package repo
    relative to the git repo.
\item pkgmod - Modify a build package.
\item pkgrel - Increment pkgrel in a package.
\item prep - Clean up a PKGBUILD file's style and find errors.
\item sitesync - Sync between a local copy of the package repo and a remote copy.
\item size-hunt - Hunt for large packages.
\item source-backup - Backup package source files.
\end{itemize}

\section{Contributing to repository}
此章告诉你如何为BlackArch Linux项目做贡献。我们接收各种PR,从小的语法修改到新的软件包。\\如果需要建议或帮助，请随时联系我们。
\\\\
欢迎所有人来贡献代码。我们感激一切贡献。

\subsection{Required tutorials}
在贡献代码之前请阅读以下内容:
\begin{itemize}
\item
\href{https://wiki.archlinux.org/index.php/Arch\_Packaging\_Standards)}{Arch
Packaging Standards}
\item \href{https://wiki.archlinux.org/index.php/Creating\_Packages}{Creating
Packages}
\item \href{https://wiki.archlinux.org/index.php/PKGBUILD}{PKGBUILD}
\item \href{https://wiki.archlinux.org/index.php/Makepkg}{Makepkg}
\end{itemize}

\subsection{Steps for contributing}
为了提交你的更改，请按照下列步骤进行:
\begin{enumerate}
\item Fork the repository from
\url{https://github.com/BlackArch/blackarch}
\item Hack the necessary files, (e.g. PKGBUILD, .patch files, etc).
\item Commit your changes.
\item Push your changes.
\item Ask us to merge in your changes, preferably through a pull request.
\end{enumerate}

\subsection{Example}
The following example demonstrates submitting a new package to the BlackArch
project. We use \href{https://wiki.archlinux.org/index.php/yaourt}{yaourt}
(you can use pacaur as well) to fetch a pre-existing PKGBUILD file for
\textbf{nfsshell} from the \href{https://aur.archlinux.org/}{AUR} and adjust it
according to our needs.

\subsubsection{Fetch PKGBUILD}
Fetch the \textit{PKGBUILD} file using yaourt or pacaur:
\begin{lstlisting}
  user@blackarchlinux $ yaourt -G nfsshell
  ==> Download nfsshell sources
  x LICENSE
  x PKGBUILD
  x gcc.patch
  user@blackarchlinux $ cd nfsshell/
\end{lstlisting}

\subsubsection{Clean up PKGBUILD}
Clean up the \textit{PKGBUILD} file and save some time:
\begin{lstlisting}
  user@blackarchlinux nfsshell $ ./blarckarch/scripts/prep PKGBUILD
  cleaning 'PKGBUILD'...
  expanding tabs...
  removing vim modeline...
  removing id comment...
  removing contributor and maintainer comments...
  squeezing extra blank lines...
  removing '|| return'...
  removing leading blank line...
  removing $pkgname...
  removing trailing whitespace...
\end{lstlisting}

\subsubsection{Adjust PKGBUILD}
Adjust the \textit{PKGBUILD} file:
\begin{lstlisting}
  user@blackarchlinux nfsshell $ vi PKGBUILD
\end{lstlisting}

\subsubsection{Build the package}
Build the package:
\begin{lstlisting}user@blackarchlinux nfsshell $ makepkg -sf
==> Making package: nfsshell 19980519-1 (Mon Dec  2 17:23:51 CET 2013)
==> Checking runtime dependencies...
==> Checking buildtime dependencies...
==> Retrieving sources...
-> Downloading nfsshell.tar.gz...
% Total    % Received % Xferd  Average Speed   Time    Time     Time
CurrentDload  Upload   Total   Spent    Left  Speed100 29213  100 29213    0
0  48150      0 --:--:-- --:--:-- --:--:-- 48206
-> Found gcc.patch
-> Found LICENSE
...
<lots of build process and compiler output here>
...
==> Leaving fakeroot environment.
==> Finished making: nfsshell 19980519-1 (Mon Dec  2 17:23:53 CET 2013)
\end{lstlisting}

\subsubsection{Install and test the package}
Install and test the package:
\begin{lstlisting}
  user@blackarchlinux nfsshell $ pacman -U nfsshell-19980519-1-x86_64.pkg.tar.xz
  user@blackarchlinux nfsshell $ nfsshell # test it
\end{lstlisting}

\subsubsection{Add, commit and push package}
Add, commit and push the package
\begin{lstlisting}user@blackarchlinux nfsshell $ cd /blackarchlinux/packages
user@blackarchlinux ~/blackarchlinux/packages $ mv ~/nfsshell .
user@blackarchlinux ~/blackarchlinux/packages $ git commit -am nfsshell && git push
\end{lstlisting}

\subsubsection{Create a pull request}
Create a pull request on \href{https://github.com/}{github.com}
\begin{lstlisting}
  firefox https://github.com/<contributor>/blackarchlinux
\end{lstlisting}

\subsubsection{Adding a remote for upstream}
A smart thing to do if you're working upstream and on a fork is to pull your own fork and add the main ba repo as a remote
\begin{lstlisting}
  user@blackarchlinux ~/blackarchlinux $ git remote -v
  origin <the url of your fork> (fetch)
  origin <the url of your fork> (push)
  user@blackarchlinux ~/blackarchlinux $ git remote add upstream https://github.com/blackarch/blackarch
  user@blackarchlinux ~/blackarchlinux $ git remote -v
  origin <the url of your fork> (fetch)
  origin <the url of your fork> (push)
  upstream https://github.com/blackarch/blackarch (fetch)
  upstream https://github.com/blackarch/blackarch (push)
\end{lstlisting}

By default, git should push straight to origin, but make sure your git config is
configured correctly. This won't be an issue unless you have commit rights as
you won't be able to push upstream without them.

If you do have the ability to commit, you might have more success using
\textit{git@github.com:blackarch/blackarch.git} but it's up to you.

\subsection{Requests}
\begin{enumerate}
\item Don't add \textbf{Maintainer} or \textbf{Contributor} comments to
\textit{PKGBUILD} files. Add maintainer and contributor names to the
AUTHORS section of BlackArch guide.
\item For the sake of consistency, please follow the general style of the other
\textit{PKGBUILD} files in the repo and use two-space indentation.
\end{enumerate}

\subsection{General tips}
\href{http://wiki.archlinux.org/index.php/Namcap}{namcap} can check packages for
errors.

%------------------%
%  Chapter 4       %
%------------------%

\chapter{Tools Guide}
敬请期待...

\section{Coming Soon}
敬请期待...

%%% APPENDIX %%%
\appendix
%%%%%%%%%%%%%%%%%%%%%%%%%%%%%%%%%%%%%%%%%%%%%%%%%%%%%%%%%%%%%%%%%%%%%%%%%%%%%%%%
%                                                                              %
% BlackArch Linux appendix template                                            %
%                                                                              %
%%%%%%%%%%%%%%%%%%%%%%%%%%%%%%%%%%%%%%%%%%%%%%%%%%%%%%%%%%%%%%%%%%%%%%%%%%%%%%%%

\appendix

\chapter{Appendix}

\section{FAQs}

\section{AUTHORS}
The following people have contributed directly to BlackArch:
\begin{itemize}
\item Tyler Bennnett (tylerb@trix2voip.com)
\item fnord0 (fnord0@riseup.net)
\item nrz (nrz@nullsecurity.net)
\item Ellis Kenyo (elken.tdos@gmail.com)
\item CaledoniaProject (the.warl0ck.1989@gmail.com)
\item sudokode (sudokode@gmail.com)
\item Valentin Churavy (v.churavy@gmail.com)
\item Boy Sandy Gladies Arriezona (reno.esper@gmail.com)
\item Mathias Nyman
\item Johannes Löthberg (demizide@gmail.com)
\item Thiago da Silva Teixeira (teixeira.zeus@gmail.com)
\end{itemize}
The following people contributed directly to ArchPwn, which has been merged
into BlackArch:
\begin{itemize}
\item Francesco Piccinno (stack.box@gmail.com)
\item jensp (jens@jenux.homelinux.org)
\item Valentin Churavy (v.churavy@gmail.com)
\end{itemize}
We have taken build code from the following people:
\begin{itemize}
\item  3ED (krzysztof1987@gmail.com)
\item  AUR Perl (aurperl@juster.info)
\item  Aaron Griffin (aaron@archlinux.org)
\item  Abakus (java5@arcor.de)
\item  Adam Wolk (netprobe@gmail.com)
\item  Aleix Pol (aleixpol@kde.org)
\item  Aleshus (aleshusi@gmail.com)
\item  Alessandro Pazzaglia (jackdroido@gmail.com)
\item  Alessandro Sagratini (ale\_sagra@hotmail.com)
\item  Alex Cartwright (alexc223@googlemail.com)
\item  Alexander De Sousa (archaur.xandy21@spamgourmet.com)
\item  Alexander Rødseth (rodseth@gmail.com)
\item  Allan McRae (allan@archlinux.org)
\item  AmaN (gabroo.punjab.da@gmail.com)
\item  Andre Klitzing (aklitzing@online.de)
\item  Andrea Scarpino (andrea@archlinux.org)
\item  Andreas Schönfelder (passtschu@freenet.de)
\item  Andrej Gelenberg (andrej.gelenberg@udo.edu)
\item  Angel Velasquez (angvp@archlinux.org)
\item  Antoine Lubineau (antoine@lubignon.info)
\item  Anton Bazhenov (anton.bazhenov@gmail.com)
\item  Arkham (arkham@archlinux.us)
\item  Arthur Danskin (arthurdanskin@gmail.com)
\item  Balda (balda@balda.ch)
\item  Balló György (ballogyor+arch@gmail.com)
\item  Bartek Piotrowski (barthalion@gmail.com)
\item  Bartosz Feński (fenio@debian.org)
\item  Bartłomiej Piotrowski (nospam@bpiotrowski.pl)
\item  Bogdan Szczurek (thebodzio@gmail.com)
\item  Brad Fanella (bradfanella@archlinux.us)
\item  Brian Bidulock (bidulock@openss7.org)
\item  C Anthony Risinger (anthony@xtfx.me)
\item  CRT (crt.011@gmail.com)
\item  Can Celasun (dcelasun@gmail.com)
\item  Chaniyth (chaniyth@yahoo.com)
\item  Chris Brannon (cmbrannon79@gmail.com)
\item  Chris Giles (Chris.G.27@gmail.com) \& daschu117
\item  Christoph Siegenthaler (csi@gmx.ch)
\item  Christoph Zeiler (archNOSPAM@moonblade.org)
\item  Clément DEMOULINS (clement@archivel.fr)
\item  Corrado Primier (bardo@aur.archlinux.org)
\item  Daenyth (Daenyth+Arch@gmail.com)
\item  Dale Blount (dale@archlinux.org)
\item  Damir Perisa (damir.perisa@bluewin.ch)
\item  Dan Fuhry (dan@fuhry.us)
\item  Dan Serban (dserban01@yahoo.com)
\item  Daniel A. Campoverde Carrión
\item  Daniel Golle
\item  Daniel Griffiths (ghost1227@archlinux.us)
\item  Daniel J Griffiths (ghost1227@archlinux.us)
\item  Daniel Micay (danielmicay@gmail.com)
\item  Dave Reisner (dreisner@archlinux.org)
\item  Dawid Wrobel (cromo@klej.net)
\item  Devaev Maxim (mdevaev@gmail.com)
\item  Devin Cofer (ranguvar@archlinux.us)
\item  DigitalPathogen (aur@InfoSecResearchLabs.co.uk)
\item  DigitalPathogen (aur@digitalpathogen.co.uk)
\item  Dmitry A. Ilyashevich (dmitry.ilyashevich@gmail.com)
\item  Dominik Heidler (dheidler@gmail.com)
\item  DrZaius (lou@fakeoutdoorsman.com)
\item  Ebubekir KARUL (ebubekirkarul@yandex.com)
\item  Eduard "bekks" Warkentin (eduard.warkentin@gmail.com)
\item  Elmo Todurov (todurov@gmail.com)
\item  Emmanuel Gil Peyrot (linkmauve@linkmauve.fr)
\item  Eric Belanger (eric@archlinux.org)
\item  Ermak (ermak@email.it)
\item  Evangelos Foutras (evangelos@foutrelis.com)
\item  Fabian Melters (melters@gmail.com)
\item  Fabiano Furtado (fusca14@gmail.com)
\item  Federico Quagliata (ntp@quaqo.org)
\item  Firmicus (francois.archlinux@org)
\item  Florian Pritz (bluewind@jabber.ccc.de)
\item  Florian Pritz (flo@xinu.at)
\item  Francesco Piccinno (stack.box@gmail.com)
\item  François Charette (francois@archlinux.org)
\item  Gaetan Bisson (bisson@archlinux.org)
\item  Geoffroy Carrier (geoffroy.carrier@koon.fr)
\item  Georg Grabler (STiAT)
\item  George Hilliard (gh403@msstate.edu)
\item  Gerardo Exequiel Pozzi (vmlinuz386@yahoo.com.ar)
\item  Gilles CHAUVIN (gcnweb@gmail.com)
\item  Giovanni Scafora (giovanni@archlinux.org)
\item  Gordin (9ordin@gmail.com)
\item  Guillaume ALAUX (guillaume@archlinux.org)
\item  Guillermo Ramos (0xwille@gmail.com)
\item  Gustavo Alvarez (sl1pkn07@gmail.com)
\item  Hugo Doria (hugo@archlinux.org)
\item  Hyacinthe Cartiaux (hyacinthe.cartiaux@free.fr)
\item  James Fryman (jfryman@gmail.com)
\item  Jan "heftig" Steffens (jan.steffens@gmail.com)
\item  Jan de Groot (jgc@archlinux.org)
\item  Jaroslav Lichtblau (dragonlord@aur.archlinux.org)
\item  Jaroslaw Swierczynski (swiergot@aur.archlinux.org)
\item  Jason Chu (jason@archlinux.org)
\item  Jason R Begley (jayray@digitalgoat.com)
\item  Jason Rodriguez
\item  Jason St. John (jstjohn@purdue.edu)
\item  Jawmare (victor2008@gmail.com)
\item  Jeff Mickey (jeff@archlinux.org)
\item  Jens Pranaitis (jens@chaox.net)
\item  Jens Pranaitis (jens@jenux.homelinux.org)
\item  Jinx (jinxware@gmail.com)
\item  John D Jones III (jnbek1972@gmail.com)
\item  John Proctor (jproctor@prium.net)
\item  Jon Bergli Heier (snakebite@jvnv.net)
\item  Jonas Heinrich
\item  Jonathan Steel (jsteel@aur.archlinux.org)
\item  Joris Steyn (jorissteyn@gmail.com)
\item  Josh VanderLinden (arch@cloudlery.com)
\item  Jozef Riha (jose1711@gmail.com)
\item  Judd Vinet (jvinet@zeroflux.org)
\item  Juergen Hoetzel (jason@archlinux.org)
\item  Juergen Hoetzel (juergen@archlinux.org)
\item  Justin Davis (jrcd83@gmail.com)
\item  Kaiting Chen (kaitocracy@gmail.com)
\item  Kaos
\item  Kevin Piche (kevin@archlinux.org)
\item  Kory Woods (kory@virlo.net)
\item  Kyle Keen (keenerd@gmail.com)
\item  Larry Hajali (larryhaja@gmail.com)
\item  LeCrayonVert
\item  Le\_suisse (lesuisse.dev+aur@gmail.com)
\item  Lekensteyn (lekensteyn@gmail.com)
\item  Limao Luo (luolimao+AUR@gmail.com)
\item  Lucien Immink
\item  Lukas Fleischer (archlinux@cryptocrack.de)
\item  Manolis Tzanidakis
\item  Marcin "avalan" Falkiewicz (avalatron@gmail.com)
\item  Mariano Verdu (verdumariano@gmail.com)
\item  Marti Raudsepp (marti@juffo.org)
\item  MatToufoutu (mattoufootu@gmail.com)
\item  Matthew Sharpe (matt.sharpe@gmail.com)
\item  Mauro Andreolini (mauro.andreolini@unimore.it)
\item  Max Pray a.k.a. Synthead (synthead@gmail.com)
\item  Max Roder (maxroder@web.de)
\item  Maxwell Pray a.k.a. Synthead (synthead@gmail.com)
\item  Maxwell Pray a.k.a. Synthead (synthead1@gmail.com)
\item  Mech (tiago.bmp@gmail.com)
\item  Michael Düll (mail@akurei.me)
\item  Michael P (ptchinster@archlinux.us)
\item  Michal Krenek (mikos@sg1.cz)
\item  Michal Zalewski (lcamtuf@coredump.cx)
\item  Miguel Paolino (mpaolino@gmail.com)
\item  Miguel Revilla (yo@miguelrevilla.com)
\item  Mike Roberts (noodlesgc@gmail.com)
\item  Mike Sampson (mike@sambodata.com)
\item  Nassim Kacha (nassim.kacha@gmail.com)
\item  Nicolas Pouillard (nicolas.pouillard@gmail.com)
\item  Nicolas Pouillard https://nicolaspouillard.fr
\item  Niklas Schmuecker
\item  Oleander Reis (oleander@oleander.cc)
\item  Olivier Le Moal (mail@olivierlemoal.fr)
\item  Olivier Médoc "oliv" (o\_medoc@yahoo.fr)
\item  Pascal E. (archlinux@hardfalcon.net)
\item  Patrick Leslie Polzer (leslie.polzer@gmx.net)
\item  Paul Mattal (paul@archlinux.org)
\item  Paul Mattal (pjmattal@elys.com)
\item  Pengyu CHEN (cpy.prefers.you@gmail.com)
\item  Peter Wu (lekensteyn@gmail.com)
\item  Philipp 'TamCore' B. (philipp@tamcore.eu)
\item  Pierre Schmitz (pierre@archlinux.de)
\item  Pranay Kanwar (pranay.kanwar@gmail.com)
\item  Pranay Kanwar (warl0ck@metaeye.org)
\item  PyroPeter (googlemail@com.abi1789)
\item  PyroPeter (googlemail.com@abi1789)
\item  Ray Rashif (schiv@archlinux.org)
\item  Remi Gacogne
\item  Renan Fernandes (renan@kauamanga.com)
\item  Richard Murri (admin@richardmurri.com)
\item  Roberto Alsina (ralsina@kde.org)
\item  Robson Peixoto (robsonpeixoto@gmail.com)
\item  Roel Blaauwgeers (roel@ttys0.nl)
\item  Rorschach (r0rschach@lavabit.com)
\item  Ruben Schuller (shiml@orgizm.net)
\item  Rudy Matela (rudy@matela.com)
\item  Ryon Sherman (ryon.sherman@gmail.com)
\item  Sabart Otto \item  Seberm (seberm@gmail.com)
\item  SakalisC (chrissakalis@gmail.com)
\item  Sam Stuewe (halosghost@archlinux.info)
\item  SanskritFritz (SanskritFritz@gmail.com)
\item  Sarah Hay (sarahhay@mb.sympatico)
\item  Sebastian Benvenuti (sebastianbenvenuti@gmail.com)
\item  Sebastian Nowicki (sebnow@gmail.com)
\item  Sebastien Duquette (ekse.0x@gmail.com)
\item  Sebastien LEDUC (sebastien@sleduc.fr)
\item  Sebastien Leduc (sebastien@sleduc.fr)
\item  Sergej Pupykin (pupykin.s+arch@gmail.com)
\item  Sergio Rubio (rubiojr@biondofu.net)
\item  Sheng Yu (magicfish1990@gmail.com)
\item  Simon Busch (morphis@gravedo.de)
\item  Simon Legner (Simon.Legner@gmail.com)
\item  Sirat18 (aur@sirat18.de)
\item  SpepS (dreamspepser@yahoo.it)
\item  Spider.007 (archPackage@spider007.net)
\item  Stefan Seering
\item  Stephane Travostino (stephane.travostino@gmail.com)
\item  Stéphane Gaudreault (stephane@archlinux.org)
\item  Sven Kauber (celeon@gmail.com)
\item  Sven Schulz (omee@archlinux.de)
\item  Sébastien Duquette (ekse.0x@gmail.com)
\item  Sébastien Luttringer (seblu@archlinux.org)
\item  TDY (tdy@gmx.com)
\item  Teemu Rytilahti (tpr@iki.fi)
\item  Testuser\_01
\item  Thanx (thanxm@gmail.com)
\item  Thayer Williams (thayer@archlinux.org)
\item  Thomas S Hatch (thatch45@gmail.com)
\item  Thorsten Töpper
\item  Tilmann Becker (tilmann.becker@freenet.de)
\item  Timothy Redaelli (timothy.redaelli@gmail.com)
\item  Timothée Ravier (tim@siosm.fr)
\item  Tino Reichardt
\item  Tobias Kieslich (tobias@justdreams.de)
\item  Tobias Powalowski (tpowa@archlinux.org)
\item  Tom K (tomk@runbox.com)
\item  Tom Newsom (Jeepster@gmx.co.uk)
\item  Tomas Lindquist Olsen (tomas@famolsen.dk)
\item  Travis Willard (travisw@wmpub.ca)
\item  Valentin Churavy (v.churavy@gmail.com)
\item  ViNS (gladiator@fastwebnet.it)
\item  Vlatko Kosturjak (kost@linux.hr)
\item  Wes Brown (wesbrown18@gmail.com)
\item  William Rea (sillywilly@gmail.com)
\item  Xavier Devlamynck (magicrhesus@ouranos.be)
\item  Xiao\item Long Chen (chenxiaolong@cxl.epac.to)
\item  aeolist (aeolist@hotmail.com)
\item  ality@pbrane.org
\item  astaroth (astaroth\_@web.de)
\item  bender02@archlinux.us
\item  billycongo (billycongo@gmail.com)
\item  bslackr (brendan@vastactive.com)
\item  cbreaker (cbreaker@tlen.pl)
\item  chimeracoder (dev@chimeracoder.net)
\item  damir (damir@archlinux.org)
\item  danitool
\item  darkapex (me@jailuthra.in)
\item  daronin
\item  dkaylor (dpkaylor@gmail.com)
\item  dobo (dobo90\_at\_gmail@com)
\item  dorphell (dorphell@archlinux.org)
\item  evr (evanroman@at.gmail)
\item  fnord0 (fnord0@riseup.net)
\item  fxbru (frxbru@gmail)
\item  hcar
\item  icarus (icarus.roaming@gmail.com)
\item  iceman (icemanf@gmail.com)
\item  kastor (kastor@fobos.org)
\item  kfgz (kfgz@interia.pl)
\item  linuxSEAT (linuxSEAT@gmail.com)
\item  m4xm4n (max@maxfierke.com)
\item  mar77i (mysatyre@gmail.com)
\item  marc0s (marc0s@fsfe.org)
\item  mickael9 (mickael9@gmail.com)
\item  nblock (nblock@archlinux.us)
\item  nofxx (x@nofxx.com)
\item  onny (onny@project
\item  pootzko (pootzko@gmail.com)
\item  revel (revel@muub.net)
\item  rich\_o (rich\_o@lavabit.com)
\item  s1gma (s1gma@mindslicer.com)
\item  sandman (r.coded@gmail.com)
\item  sebikul (sebikul@gmail.com)
\item  sh0 (mee@sh0.org)
\item  shild (sxp@bk.ru)
\item  simo (simo@archlinux.org)
\item  snuo
\item  sudokode (sudokode@gmail.com)
\item  tobias (tobias@archlinux.org)
\item  trashstar (trash@ps3zone.org)
\item  unexist (unexist@subforge.org)
\item  untitled (rnd0x00@gmail.com)
\item  virtuemood (virtue@archlinux.us)
\item  wido (widomaker2k7@gmail.com)
\item  wodim (neikokz@gmail.com)
\item  yannsen (ynnsen@gmail.com)
\end{itemize}


\end{document}

%%% EOF %%%
